\documentclass[../main.tex]{subfiles}
\begin{document}
%%%%%%%%%%%%%%%%%%%%%%%%%%%%%%%%%%%%%%%%%%%%%%%%%%%%%%%%%%%%%%%%%%%%%%%%%%%%%%%%%%%%%%%%%%%%%%%%%%%
\chapter{Using Artboards (\texorpdfstring{\textbf{PRO}}{PRO})}

As you learned in the previous chapter, your OmniGraffle projects can be as simple or complex as
you need. For example, your project can have a single canvas with a single layer, a single
canvas with multiple layers, or it can contain multiple canvases each with multiple layers.

This chapter introduces you to a new layer---and object---type that's only available in
OmniGraffle Pro: The \textbf{Artboard}. \textbf{Artboard Layers} reside beneath or between the
layers of your project, and the objects on an Artboard Layer, known as \textbf{Artboard
Objects}, have a direct effect on the objects residing on layers higher up in the stack.
Artboards can be used to define export areas in your project, while also acting as a container
for objects on the layers above.

For example, you could use artboards to:

\begin{itemize}
  \item Define the areas needed for toolbar buttons in a macOS app. You could mock up the
    toolbar, and use separate artboards for each button location.
  \item Export multiple resolution images (normal, 100\%, through @5x, 500\%) for use in an app
    bundle or for a website.
  \item Act as containers for objects in a project, such as an app prototype, where you need the
    ability to move interaction ``screens'' around without breaking everything else.
  \item And much, much more\ldots{}
\end{itemize}

As you'll soon discover, artboards are powerful new tools for your bag of design tricks.

\section{Understanding Artboards}

The fastest way to add an Artboard Layer \textit{and} an Artboard Object to your project is with
the \textbf{Artboard Tool}:

\ogfigure{0.800}{res/og7_artboards_toolpalette.png}

Artboards are unlike other layers in OmniGraffle. For example, if you create a new OmniGraffle
project and add a \textbf{New Layer}, you'll notice that \textbf{Layer 2} appears
\textit{above} \textbf{Layer 1} in the layer stack.

\ogfigure{0.800}{res/og7_artboards01.png}

However, with the same new project, if you choose \textbf{File} $\blacktriangleright$
\textbf{New Artboard Layer} in the toolbar (or use the Artboard Tool to create a new artboard
object and layer), you'll notice that the \textbf{Artboard Layer 1} layer is added to the
canvas \textit{underneath} \textbf{Layer 1}.

\ogfigure{0.800}{res/og7_artboards02.png}

This subtle behavior is deliberate because artboards are used to manage the behavior and output
of the things you create on the layers \textit{above}. In fact, whenever you add a new
Artboard Layer, it pops into the layer stack underneath the selected layer.

And, like standard layers, your project can have as many Artboard Layers that you need; but
unlike standard layers, there is no requirement for a project to have at least one Artboard
Layer. If you don't need one, don't use one. (The same applies to Shared Layers, too.)

\subsection{How Artboard Objects Interact with Objects on Other Layers}

As noted earlier, artboards interact with the objects on layers higher up in the stack. This
means that when you go to grab and move an artboard to a different location on the canvas,
every object that's on a higher layer in the project---and within the bounds of the artboard,
partially or otherwise---goes along for the ride.

To help illustrate the point, let's look at a simple example\ldots{}

\begin{enumerate}
  \item Create a new OmniGraffle project.
  \item In the menu bar, choose \textbf{View} and look in the \textbf{Extras} section to see if
    \textbf{Movement Handles} has a checkmark next to it.
  \item Click the \textbf{New Layer} button in the toolbar to add another standard layer to the
    project; don't worry about renaming this layer.
  \item On \textbf{Layer 1}, use the Shape tool to draw a square.
  \item On \textbf{Layer 2}, use the Shape tool to draw a circle, but have the circle overlap
    the square a little, like so:

    \ogfigure{0.800}{res/og7_artboards_interaction01.png}
  \item In the Tool Palette, select the \textbf{Artboard Tool}
    \ogiconw{19.4pt}{res/og7_tools_artboard.png}.
  \item Draw an artboard so that it intersects the circle on Layer 2:

    \ogfigure{0.800}{res/og7_artboards_interaction02.png}

    If you look at the Sidebar, you'll notice that the Artboard Layer has been added between
    Layers 1 and 2.
  \item Drag the Artboard Layer to the bottom of the layer stack, so that Layers 1 and 2 appear
    above the artboard.

    \ogfigure{0.673}{res/og7_artboards_interaction03.png}
\end{enumerate}

With the artboard at the bottom of the layer stack, and knowing that artboard objects affect the
objects on layers higher up in the stack, \textbf{click and drag} on the artboard. Notice how
the circle on Layer 2 moves along with the artboard but the square on Layer 1 stays in place.
This is because a portion of the circle overlaps the artboard, but the square does not.

\ogfigure{0.673}{res/og7_artboards_interaction04.png}

Press \textbf{Command-Z} to undo that movement and reset the circle and artboard to their
previous locations.

If all you wanted to do was move the artboard so that it is placed beneath the square
\textit{and} the circle, here's how you would do that:

\begin{enumerate}
  \item In the Sidebar, select the artboard.
  \item First, \textbf{click-and-hold} on the artboard's Movement Handle, located at the center
    of the Bounding Box Basics. This grabs the artboard so it can be moved:

    \ogfigure{0.673}{res/og7_artboards_interaction05.png}
  \item \textit{Then} press the \textbf{Control} key and move the artboard into position:

    \ogfigure{0.673}{res/og7_artboards_interaction06.png}

    When the Control key is added in that way, it temporarily decouples the artboard from any
    objects on layers higher up in the stack so you can freely move the artboard around.
  \item When the artboard is in the correct location, let go of the mouse button and then let go
    of the Control key.
\end{enumerate}

Now if you move the artboard, both the square and circle go along for the ride.

\subsection{Using Artboards as Containers}

Sometimes, you need the flexibility of a group, without actually creating a group. For example,
when you're trying to rapidly prototype something and you want the flexibility to move things
around to see how they look when you move or add things to the design.

For example, let's say that you're working on toolbar button designs for a Mac app. Toolbars, as
you know, contain three circular window control buttons in the upper-left corner. One quick
way to attach these to an artboard is to first select the three buttons, and then choose
\textbf{Arrange} $\blacktriangleright$ \textbf{Make Artboard} (\textbf{Shift-Command-A}).

\ogfigure{0.800}{res/og7_artboards22.png}

The \textbf{Make Artboard} option wraps the selected objects in an artboard, and places that
artboard layer at the bottom of the stack.

\ogfigure{0.800}{res/og7_artboards22a.png}

If, while iterating on your design, you decide that you want to see how the toolbar looks
without the window controls, you can just grab the artboard and move it off to the side.

\ogfigure{0.673}{res/og7_artboards23.png}

The objects are still visible on the canvas and can be dragged back into place at any time.

\subsection{Adding a Toolbar Button for Creating New Artboard Layers}

If you plan to work with Artboards often, you should consider adding the \textbf{New Artboard
Layer} button to OmniGraffle's toolbar. To do this, follow these steps:

\begin{enumerate}
  \item \textbf{Control-click} on the toolbar, and choose \textbf{Customize Toolbar}
  \item Locate the \textbf{New Artboard Layer} button and drag that up to the toolbar:

    \ogfigure{0.165}{res/og7_button_newartboardlayer.png}

    \ogfigure{0.800}{res/og7_artboards_toolbar.png}
  \item Click \textbf{Done}
\end{enumerate}

If you ever decide that you no longer want the New Artboard Layer button in the toolbar,
\textbf{Control-click} on the toolbar again, and then drag the \textbf{New Artboard Layer}
button away from the toolbar. When you let go of the mouse button, it disappears in a poof.

\subsection{Changing an Artboard Object's Shape}

By default, the Artboard Tool draws a rectangular shape when used. However, you can use the
\textbf{Shape} inspector to change the shape of an Artboard Object \textit{before} or
\textit{after} it has been added to an Artboard Layer; here's how:

\subsubsection{Before Drawing the Artboard Object}

\begin{enumerate}
  \item In the Tool Palette, click to select the \textbf{Artboard Tool}
  \item Press \textbf{Command-1} to view the \textbf{Object} inspectors in the inspector bar
  \item \textbf{Option-click} the \textbf{Shape} inspector; this opens the Shape inspector, and
    collapses the other Object inspectors
  \item Click the \textbf{Shape Selection} button at the top of the \textbf{Shape} inspector,
    and then double-click on another shape in the pop-up menu that appears.
\end{enumerate}

\subsubsection{After Drawing the Artboard Object}

\begin{enumerate}
  \item Select the Artboard Object on the Canvas
  \item \textbf{Option-click} the \textbf{Shape} inspector; this opens the Shape inspector, and
    collapses the other Object inspectors
  \item Click the \textbf{Shape Selection} button at the top of the \textbf{Shape} inspector,
    and then click on one of the shapes in the pop-up menu that appears. The Artboard Object
    instantly changes its shape.
\end{enumerate}

\subsection{Deleting Artboards}

If you decide that your project no longer needs an artboard, you can safely delete the artboard
without deleting the objects that reside on other layers.

To delete an artboard, you can do one of the following:

\begin{itemize}
  \item \textbf{Control-click} on the artboard layer in the Sidebar and choose \textbf{Delete
    Layer} from the contextual menu that appears.
  \item Select the artboard object in the Sidebar or on the canvas and press \textbf{Delete}.
\end{itemize}

\begin{ogwarning}
  If you select the artboard on the canvas, make sure that you select \textbf{just} the artboard,
  and not any of the objects on the layers above prior to pressing \textbf{Delete}.

  If you delete something by accident, remember that you can always press \textbf{Command-Z} to
  undo that action.
\end{ogwarning}

\section{Working with Artboards}

To help illustrate how to use artboards in your own projects, there are a couple approaches you
can take. First, if you know the dimensions of the image assets you need, you can start by
laying down your artboard first and then drawing on the layers above. The second approach is
when you already have the art you need in OmniGraffle, and you want to add an artboard to
define the export areas you need.

\subsection{Defining the Design Area with Artboards}

Let's say that you've started up your own private rocket company, \textbf{MasterBlasters}, and
you're using OmniGraffle for all of your design needs---from designing your brand ID to
prototyping your website and companion app for iOS.

The Web Team needs a transparent PNG image that's 1080 $\times$ 256 pixels, and the App Team
needs a similar PNG image, except it needs to be 512 $\times$ 512 pixels.

In the past, the process for doing this would be extremely tedious. You'd have different project
files with different-sized assets, which meant there was a greater chance for failure when you
split up everything that way---and who wants that when you're creating a rocket company?

However, to accomplish this feat, you'll only need two artboard objects. To get the larger
graphic for the billboard, you'll repurpose the artboard object you create to define the web
banner, and then change the \textbf{Unit Scale} in the \textbf{Units} inspector.

This is where the \textit{power} of artboard layers and objects come into play. Let's get
started!

\begin{enumerate}
  \item Create a new OmniGraffle project.
  \item Select the \textbf{Artboard Tool}, hold down the \textbf{Shift} key, and then draw a
    \textbf{256 $\times$ 256 px} square on the canvas:

    \ogfigure{0.800}{res/og7_artboards_area01.png}

    Holding down the \textbf{Shift} key while drawing a shape ensures that the height and width of
    the object are the same dimensions.
  \item \textbf{Option-click} on the \textbf{Shape} inspector and set the Corner Radius to
    \textbf{8 pt}:

    \ogfigure{0.800}{res/og7_artboards_area02.png}
  \item In the Sidebar, double-click on the artboard's name (currently, \textbf{Rectangle}), and
    change that to \textbf{app-icon}:

    \ogfigure{0.800}{res/og7_artboards_area03.png}

    Notice how the Artboard Layer has a red highlight to its layer preview in the Sidebar. This
    makes it easy for you to quickly identify the artboards in your project.

    \begin{ogtip}
      If you are designing iOS app icons, there is an \textbf{iOS App Icon} shape in the
      \textbf{Shape} inspector's Shape Popover.
    \end{ogtip}
  \item Click to select the artboard on the canvas.
  \item Hold down the \textbf{Option} key, and then \textbf{click and drag} on the artboard to
    create a duplicate:

    \ogfigure{0.673}{res/og7_artboards_area04.png}
  \item In the Sidebar, double-click on the duplicate artboard's name, and change that to
    \textbf{web-banner}:

    \ogfigure{0.673}{res/og7_artboards_area05.png}
  \item Click to select the \textbf{web-banner} artboard on the canvas.
  \item \textbf{Option-click} the \textbf{Geometry} inspector, and change the artboard's width
    to \textbf{1080} pixels:

    \ogfigure{0.800}{res/og7_artboards_area06.png}
  \item Next, you'll want to align the artboards so that the \textbf{app-icon} artboard is
    centered on the \textbf{web-banner} artboard. Click on the canvas to bring that area into
    focus, and then press \textbf{Command-A} (\textbf{Select All}) to select both artboards.
  \item \textbf{Option-click} the \textbf{Alignment} inspector, and then press the following
    buttons:

    \begin{itemize}
      \item \ogicon{12.5pt}{res/og7_button_align_vertcenter.png} to align the artboards based on
        their vertical centers
      \item \ogicon{12.5pt}{res/og7_button_align_horzcenter.png} to align the artboards based on
        their horizontal centers
    \end{itemize}
\end{enumerate}

This is how the artboards should appear on the canvas:

\ogfigure{0.800}{res/og7_artboards_area07.png}

With the artboards in place, you can work on your design on the layers above. When you set up
your OmniGraffle project, you layer it just like you normally would, with a background, the
rocketship, and the text blocks you need for the web banner.

\ogfigure{0.800}{res/og7_artboards_area08.png}

\subsubsection{Exporting Image Assets with Artboards}

To create the image files you need, choose \textbf{File} $\blacktriangleright$ \textbf{Export}
(\textbf{Option-Command-E}). After choosing \textbf{Export}, you'll catch your first glimpse
of OmniGraffle's new Export panel.

\ogfigure{0.800}{res/og7_artboards13.png}

On the left side of the Export panel, under the \textbf{Export Area} section, make sure that
\textbf{Export from artboards only} is checked. This tells OmniGraffle to export images based
on the artboards in your project; without this option checked, what gets exported depends on
what you have selected in the Export Area pop-up menu.

If you look at the right side of the Export panel, you'll see the \textbf{Export Preview}, where
you can page through the images as determined by their artboards.

\ogfigure{0.800}{res/og7_artboards14.png}

You'll notice that, when exported at 100\% size, the \textbf{web-banner} image will export at
the size you need (1080 $\times$ 256 pixels), while the \textbf{app-icon} image exports at the
size of the artboard---256 $\times$ 256 pixels---except, you need the \textbf{app-icon} to be
512 $\times$ 512 pixels.

If you look at the \textbf{Export Sizes} area of the Export panel, you'll see that you can also
specify variable export sizes:

\ogfigure{0.673}{res/og7_artboards_export01.png}

In the \textbf{Export Sizes} area, click \ogicon{12.5pt}{res/og7_button_export_addsize.png} to
add another export size. By default, OmniGraffle sets the next size to \textbf{200\%} and adds
\textbf{@2x} as the suffix.

\ogfigure{0.673}{res/og7_artboards_export02.png}

If you look at the Export Preview, you'll notice that there are now four images that will be
exported: two at 100\% and two at 200\%:

\ogfigure{0.673}{res/og7_artboards_export03.png}

When you click \textbf{Export}, the Save sheet flops down so you can choose a location on your
Mac to save the files. When you export, the text you enter in the \textbf{Export As} field is
used as the name for a folder to contain the images.

\ogfigure{0.800}{res/og7_artboards15.png}

When you click \textbf{Save}, OmniGraffle outputs the image files to the folder you have
specified. In this case, four images are exported:

\begin{itemize}
  \item app-icon.png (256 $\times$ 256 px)
  \item app-icon@2x.png (512 $\times$ 512 px)
  \item web-banner.png (1080 $\times$ 256 px)
  \item web-banner@2x.png (2160 $\times$ 512 px)
\end{itemize}

\ogfigure{0.800}{res/og7_artboards16.png}

Depending on the needs of the Web and App Teams, you can either send them both sizes, or just
the ones they need and then discard the others.

\subsection{Adding Artboards to an Existing Project}

Now that you know what artboards look like and how they interact with objects on other layers,
you can also add artboards to an existing OmniGraffle project.

To help illustrate the point, let's use an example of a spaceship landing game. The app is
designed to fit an iPhone 7 screen (750 $\times$ 1334 pixels), both in portrait and landscape
orientation. To start, the OmniGraffle project has three layers:

\begin{itemize}
  \item One that contains the spaceship.
  \item A background layer that depicts an alien atmosphere.
  \item A mountainous foreground with a few ``designated'' landing areas among the hazardous
    peaks.
\end{itemize}

To add artboards to this project, do the following:

\begin{enumerate}
  \item Since you want the artboards to appear beneath everything else in the project, select
    the layer lowest in the stack.
  \item In the Tool Palette, select the \textbf{Artboard Tool} and then draw a \textbf{750
    $\times$ 1334 px} rectangular artboard; these are the dimensions for the Portrait display.
    When you draw the artboard, OmniGraffle creates the artboard layer to contain the artboard.
  \item In the sidebar, name the artboard layer \textbf{Portrait} and the artboard itself
    \textbf{portrait}.

    \ogfigure{0.800}{res/og7_artboards_adding01.png}
  \item Next, you need to create a \textbf{Landscape} version of the artboard, and for this,
    you'll employ a little trick. \textbf{Control-click} on the \textbf{Portrait} artboard layer
    and choose \textbf{Duplicate Layer} from the contextual menu that appears:

    \ogfigure{0.673}{res/og7_artboards_adding02.png}
  \item Rename the duplicated artboard layer to \textbf{Landscape}, and rename the artboard to
    \textbf{landscape}, and then select the \textbf{landscape} artboard in the Sidebar.

    \ogfigure{0.673}{res/og7_artboards_adding03.png}
  \item \textbf{Option-click} on the \textbf{Geometry} inspector, and change the dimensions of
    the \textbf{landscape} artboard to \textbf{1334 px} wide by \textbf{750 px} high:

    \ogfigure{0.800}{res/og7_artboards_adding04.png}
  \item After changing the dimensions of the \textbf{landscape} artboard, you'll notice that it
    needs to be repositioned:

    \ogfigure{0.800}{res/og7_artboards_adding05.png}
  \item Use the following to relocate the \textbf{landscape} artboard without moving the objects
    on the layers above:

    \begin{enumerate}
      \item \textbf{Click-and-hold} on the \textbf{landscape} artboard's center handle. This grabs
        the artboard so it can be moved.
      \item Press the \textbf{Control} key, and \textit{then} move the artboard into position:
    \end{enumerate}

    \ogfigure{0.800}{res/og7_artboards_adding06.png}

    When the \textbf{landscape} artboard is in place, let go of the mouse button and the Control
    key.
\end{enumerate}

And here are the Landscape and Portrait views of the game prototype after exporting the images:

\ogfigure{0.800}{res/og7_artboards_adding07.png}
%%%%%%%%%%%%%%%%%%%%%%%%%%%%%%%%%%%%%%%%%%%%%%%%%%%%%%%%%%%%%%%%%%%%%%%%%%%%%%%%%%%%%%%%%%%%%%%%%%%
\end{document}
